\documentclass[11pt]{article}

\usepackage[margin=1in]{geometry}
\usepackage{amsmath,amssymb,amsthm}
\usepackage{enumitem}
\usepackage{graphicx}
\usepackage[colorlinks=true,linkcolor=blue,urlcolor=blue,citecolor=blue]{hyperref}
\emergencystretch=2em

\newtheorem{theorem}{Theorem}
\newtheorem{lemma}[theorem]{Lemma}
\newtheorem{corollary}[theorem]{Corollary}
\theoremstyle{remark}
\newtheorem{remark}[theorem]{Remark}

\newcommand{\cN}{\mathcal N}
\newcommand{\covN}{\operatorname{cov}(\mathcal N)}
\newcommand{\cNM}{\mathcal{NM}}
\newcommand{\cNF}{\mathcal{NF}}
\newcommand{\cND}{\mathcal{ND}}
\newcommand{\cAN}{\mathcal{AN}}
\newcommand{\cANM}{\mathcal{ANM}}
\newcommand{\esssup}{\operatorname*{ess\,sup}}
\newcommand{\essinf}{\operatorname*{ess\,inf}}
\newcommand{\finc}{\mathfrak c}

\title{The Sharp \(\operatorname{cov}(\mathcal N)\) Threshold\\
for Three Families of Pathological Nets}
\author{Math discovery candidate partial result}
\date{August 17, 2026}

\begin{document}
\maketitle

\section{Status and source questions}

This packet gives a candidate substantial partial answer to Questions Q1--Q3
in D.~L. Rodr\'iguez-Vidanes,
\emph{Lineability on nets and uncountable sequences of functions in measure
theory}, arXiv:2306.\allowbreak03862v1 \cite{source}.  The questions occur in Section 4,
printed page 26:

\begin{center}
\includegraphics[width=\textwidth]{figures/open_problem_crop.png}
\end{center}

The result completely determines the smallest-cardinality clause of Q1 and
Q2 for every dimension requested there, and of Q3 up to dimension
\(\finc\).  It does not classify all directed sets.

\section{Definitions}

Write \(\cN\) for the ideal of Lebesgue-null subsets of \([0,1]\), and
\(\covN\) for the least cardinality of a subfamily of \(\cN\) covering
\([0,1]\).  Let \(\mathcal A\) be a directed set.  We use the source paper's
families:
\begin{itemize}[leftmargin=*]
\item \(\cNM_{\mathcal A}\) consists of a.e.-increasing nets
  \((f_a)_{a\in\mathcal A}\) of nonnegative measurable functions which
  converge pointwise a.e. to an integrable \(f\), but for which monotone
  convergence of the integrals fails.
\item \(\cNF_{\mathcal A}\) consists of nets of a.e.-nonnegative measurable
  functions for which
  \[
     \int \liminf_{a\in\mathcal A} f_a
       > \liminf_{a\in\mathcal A}\int f_a,
  \]
  with the pointwise net liminf integrable.
\item \(\cND_{\mathcal A}\) consists of nets \((f_a)\), dominated a.e. at
  every index by one integrable function, which converge pointwise a.e. to
  an integrable \(f\), but for which
  \(\int|f_a-f|\not\to0\).
\end{itemize}
All integrals below are over \([0,1]\).

\section{Proof intuition}

The usual sequence proofs conceal a countability step.  In the Lebesgue
measure algebra, even an uncountable essential union or essential supremum is
determined by a countable subfamily.  If a dominated net has large
\(L^1\)-error arbitrarily far out, every index tail therefore has a
positive-measure essential union of large-error sets, and the essential
tail-limsup still has positive measure.  Pointwise convergence forces each
point of that limsup into a null exceptional set attached to some index.
Thus the index set supplies a null cover of a positive-measure set, requiring
at least \(\covN\) indices.

For monotone nets, compare the literal pointwise limit with the essential
supremum of all terms.  For Fatou, compare the literal tail infima with the
essential tail infima.  The same null-cover obstruction makes these objects
agree in the direction needed for the classical integral inequalities.

\section{Measure-algebra lemmas}

\begin{lemma}[Positive measurable sets have covering number \(\covN\)]
\label{lem:positive-cover}
If \(E\subset[0,1]\) is measurable and \(\lambda(E)>0\), and
\(E\subset\bigcup_{\xi<\kappa}N_\xi\) with every \(N_\xi\in\cN\), then
\(\kappa\ge\covN\).
\end{lemma}

\begin{proof}
Work modulo \(1\) on the circle.  There are countably many translates of
\(E\) whose union is conull.  Indeed, if \(R\) is the uncovered measurable
remainder, averaging over translations \(t\) gives
\[
 \int \lambda\bigl(R\cap(E+t)\bigr)\,dt
       =\lambda(R)\lambda(E).
\]
Choose a translate covering at least this proportion of \(R\), and iterate;
the remainder decreases geometrically to a null set.

If \(\kappa<\covN\), each selected translate of \(E\) is covered by
\(\kappa\) translates of the \(N_\xi\)'s.  These null sets, together with
the null remainder, form a cover of the circle by at most
\(\kappa\cdot\aleph_0<\covN\) null sets, contradicting the definition of
\(\covN\).
\end{proof}

\begin{lemma}[Countable determination]
\label{lem:countable}
Every essential union or essential intersection of an arbitrary family of
Lebesgue-measurable sets is equal modulo null sets to the corresponding
operation on a countable subfamily.  Consequently, every essential supremum
or essential infimum of an arbitrary family of measurable extended-real
functions is determined a.e. by a countable subfamily.
\end{lemma}

\begin{proof}
For a family \(\mathcal E\) of measurable sets, let
\[
 s=\sup\left\{\lambda\left(\bigcup\mathcal F\right):
                \mathcal F\subset\mathcal E\text{ finite}\right\}.
\]
Choose finite \(\mathcal F_n\subset\mathcal E\) for which the displayed
measures tend to \(s\), and let
\(U=\bigcup_n\bigcup\mathcal F_n\).  Then \(\lambda(U)=s\).  For every
\(E\in\mathcal E\), one has \(\lambda(E\setminus U)=0\), since otherwise
\(\lambda(U\cup E)>s\).  Hence \(U\) is the essential union.  Complements
give essential intersections.

For functions, apply the set statement to the rational superlevel sets
\(\{h_i>q\}\), \(q\in\mathbb Q\), and take the union of the countably many
resulting countable index sets.  Essential infima follow by reversing the
inequalities.
\end{proof}

\begin{lemma}[A dominated bad net needs \(\covN\) indices]
\label{lem:dominated}
Let \((h_a)_{a\in\mathcal A}\) be measurable functions such that
\(0\le h_a\le G\) a.e. for every \(a\), where \(G\in L^1[0,1]\), and suppose
\(h_a\to0\) pointwise a.e.  If \(\int h_a\not\to0\), then
\[
          |\mathcal A|\ge\covN.
\]
\end{lemma}

\begin{proof}
There is an \(\varepsilon>0\) such that every index tail contains a
\(b\) with \(\int h_b\ge\varepsilon\).  Choose \(M>0\) so that
\(\int_{\{G>M\}}G<\varepsilon/4\), put
\(\delta=\varepsilon/4\), and, for each such large \(b\), set
\[
 E_b=\{x:G(x)\le M,\ h_b(x)\ge\delta\}.
\]
The estimate
\[
 \varepsilon\le\int h_b
 \le \frac{\varepsilon}{4}+\delta+M\lambda(E_b)
\]
shows that \(\lambda(E_b)\ge\eta:=\varepsilon/(2M)>0\).

For \(a\in\mathcal A\), let \(U_a\) be the essential union of all \(E_b\)
with \(b\ge a\) and \(\int h_b\ge\varepsilon\).  Every tail contains such a
\(b\), so \(\lambda(U_a)\ge\eta\).  The family \((U_a)\) is
down-directed.  Let
\[
       U=\essinf_{a\in\mathcal A}U_a.
\]
By Lemma~\ref{lem:countable}, a countable subfamily \(U_{a_n}\) determines
\(U\).  Recursively choose \(c_n\) above \(c_{n-1}\) and \(a_n\).
Then \(U_{c_n}\) decreases to \(U\) modulo null sets, and
\[
       \lambda(U)=\lim_n\lambda(U_{c_n})\ge\eta.
\]

Again by Lemma~\ref{lem:countable}, for each \(a\) choose a countable set
\(C_a\) of large indices \(b\ge a\) such that
\[
       U_a=\bigcup_{b\in C_a}E_b\quad\text{modulo null sets}.
\]
Thus
\[
       D_a:=U\setminus\bigcup_{b\in C_a}E_b
\]
is null.  Outside the fixed null set where \(h_a\not\to0\), every
\(x\in U\) has an index \(a_x\) after which \(h_b(x)<\delta\).
Consequently \(x\notin E_b\) for \(b\in C_{a_x}\), so \(x\in D_{a_x}\).
We have covered the positive-measure set \(U\) by one fixed null set and
the family \((D_a)_{a\in\mathcal A}\).  Lemma~\ref{lem:positive-cover}
gives \(|\mathcal A|\ge\covN\).
\end{proof}

\section{The sharp lower bound}

\begin{theorem}
\label{thm:empty}
If \(\mathcal A\) is a directed set with
\(|\mathcal A|<\covN\), then
\[
       \cNM_{\mathcal A}
       =\cNF_{\mathcal A}
       =\cND_{\mathcal A}
       =\varnothing.
\]
\end{theorem}

\begin{proof}
\textbf{The dominated family.}
If \((f_a)\in\cND_{\mathcal A}\) has pointwise a.e. limit \(f\) and
\(|f_a|\le g\) a.e., apply Lemma~\ref{lem:dominated} to
\[
       h_a=|f_a-f|,\qquad G=g+|f|.
\]
This contradicts \(|\mathcal A|<\covN\).

\medskip
\noindent
\textbf{The monotone family.}
Suppose \((f_a)\in\cNM_{\mathcal A}\) has pointwise a.e. limit \(f\).
Set
\[
                    u=\esssup_{a\in\mathcal A}f_a.
\]
Choose \(a_n\) determining this essential supremum and recursively choose
\(c_n\) above \(c_{n-1}\) and \(a_n\).  The a.e. order makes
\((f_{c_n})\) increasing modulo null sets, with supremum \(u\).  Ordinary
monotone convergence yields
\[
       \int u=\sup_n\int f_{c_n}
             =\sup_{a\in\mathcal A}\int f_a.
\tag{1}
\]

We claim that \(f=u\) a.e.  Since \(\{f_a>u\}\) is null for every \(a\),
pointwise convergence gives
\[
 \{f>u\}\subset N\cup\bigcup_{a\in\mathcal A}\{f_a>u\},
\tag{2}
\]
where \(N\) is the null convergence-exception set.  Lemma
\ref{lem:positive-cover} and \(|\mathcal A|<\covN\) imply
\(\lambda\{f>u\}=0\).

For the reverse inequality, fix \(n\).  If
\(x\in\{f_{c_n}>f\}\setminus N\), convergence supplies an index after which
the values are strictly below \(f_{c_n}(x)\).  Taking a common upper bound
with \(c_n\) shows
\[
 \{f_{c_n}>f\}\subset
 N\cup\bigcup_{b\ge c_n}\{f_{c_n}>f_b\}.
\tag{3}
\]
Every set in the last union is null by the a.e. monotonicity.  The same
covering argument makes \(\{f_{c_n}>f\}\) null.  Taking the countable
supremum gives \(u\le f\) a.e.  Thus \(u=f\) a.e., and (1) is precisely
the monotone-convergence conclusion, a contradiction.

\medskip
\noindent
\textbf{The Fatou family.}
Let \((f_a)\) be a net of a.e.-nonnegative measurable functions.  Define
the literal and essential tail infima by
\[
 q_a(x)=\inf_{b\ge a}f_b(x),\qquad
 g_a=\essinf_{b\ge a}f_b,
\]
and put
\[
 f=\sup_{a\in\mathcal A}q_a
      =\liminf_{a\in\mathcal A}f_a,\qquad
 g=\esssup_{a\in\mathcal A}g_a.
\]
Assume, as in the definition of \(\cNF_{\mathcal A}\), that \(f\) is
integrable.

For each \(a\), a countable tail subfamily determines \(g_a\), so
\(q_a\le g_a\le g\) outside a null set \(D_a\).  Therefore
\[
       \{f>g\}\subset\bigcup_{a\in\mathcal A}D_a.
\]
The left side is measurable; Lemma~\ref{lem:positive-cover} implies that it
is null.  Hence \(f\le g\) a.e.

The net \((g_a)\) is increasing in the directed order.  Countably determine
its essential supremum and take successive common upper bounds as before.
Ordinary monotone convergence then gives
\[
\begin{aligned}
 \int f
 &\le \int g
  = \sup_{a\in\mathcal A}\int g_a                                      \\
 &\le \sup_{a\in\mathcal A}\inf_{b\ge a}\int f_b
  = \liminf_{a\in\mathcal A}\int f_a.
\end{aligned}
\]
This is Fatou's inequality and contradicts membership in
\(\cNF_{\mathcal A}\).
\end{proof}

\section{Consequences for Q1--Q3}

\begin{corollary}[Exact minimum in Q1 and Q2]
\label{cor:q12}
For every cardinal \(1\le\mu\le2^{\finc}\), the smallest possible
cardinality of a directed set in Q1 is
\[
                         \boxed{\covN}.
\]
The same is true for Q2.
\end{corollary}

\begin{proof}
The lower bound follows from Theorem~\ref{thm:empty}, since even one nonzero
member of any requested cone would make the relevant family nonempty.

For the upper bound, Theorem 2.4 of the source paper \cite{source} constructs
in ZFC a directed set of cardinality \(\covN\) and a positive cone of
dimension \(2^{\finc}\) contained simultaneously in
\(\cNM_{\mathcal A}\), \(\cNF_{\mathcal A}\), and
\(\cND_{\mathcal A}\), after adjoining zero.  Selecting \(\mu\) generators
gives every smaller requested dimension.  For Q2, note additionally that
\(\cND_{\mathcal A}\) is symmetric under multiplication by \(-1\), so the
same generators also give the negative cone required by the source's
definition of coneability.
\end{proof}

\begin{corollary}[Exact minimum in Q3 through \(\finc\)]
\label{cor:q3}
For every cardinal \(1\le\mu\le\finc\), the smallest possible cardinality
in Q3 is
\[
                         \boxed{\covN}.
\]
This remains true if Q3 is read as requiring one common directed set for all
three families.
\end{corollary}

\begin{proof}
The family \(\cND_{\mathcal A}\) occurring in Q3 forces the lower bound by
Theorem~\ref{thm:empty}.

For the upper bound, the source constructs directed sets of cardinality
\(\covN\) on which, respectively,
\(\cAN_{\mathcal A}\) is strongly \(2^{\finc}\)-algebrable,
\(\cND_{\mathcal A}\) is strongly \(\finc\)-algebrable, and
\(\cANM_{\mathcal A}\setminus\cND_{\mathcal A}\) is strongly
\(\finc\)-algebrable; see its Theorems 2.17, 2.13, and 2.8
\cite{source}.  Let these index sets be
\(\mathcal A_1,\mathcal A_2,\mathcal A_3\), and give
\[
              \mathcal A
              =\mathcal A_1\times\mathcal A_2\times\mathcal A_3
\]
the coordinatewise order.  Its cardinality is still \(\covN\).

Pull a net on \(\mathcal A_i\) back along the \(i\)-th coordinate
projection.  That projection is order preserving and cofinal, so the
pointwise limit and every relevant eventual integral or convergence
property are unchanged.  The pullback is an injective algebra
homomorphism, hence preserves freeness.  Thus the three source algebras live
on one common \(\mathcal A\).  Restricting their free generating sets gives
every \(1\le\mu\le\finc\).
\end{proof}

\section{Verification, novelty, and limitations}

The proof is structural and uses no numerical computation.  It was audited
for the two points most likely to hide a countability assumption:
Lemma~\ref{lem:countable} explicitly reduces essential suprema and infima to
sequences, while Lemma~\ref{lem:dominated} never assumes a countable cofinal
subset of \(\mathcal A\).  Equations (2) and (3) keep the varying a.e.
exception sets visible rather than silently choosing common
representatives.

A bounded novelty search through August 17, 2026 used the exact title,
arXiv id, author, the phrase \emph{smallest possible cardinality}, and
combinations of \(\covN\), nets, dominated convergence, monotone convergence,
and lineability.  It located the 2025 journal publication of the source and
a later paper by Dantas--Rodr\'iguez-Vidanes on different net pathologies,
but no source stating Theorem~\ref{thm:empty} or Corollaries
\ref{cor:q12}--\ref{cor:q3}.  Novelty is therefore plausible, not certified.

The result leaves open:
\begin{itemize}[leftmargin=*]
\item the classification of all directed sets in Q1--Q3;
\item Q3 for \(\finc<\mu\le2^{\finc}\);
\item the regular-cardinal Questions Q4--Q6.
\end{itemize}
The first remaining issue depends on the cofinal/Tukey structure of the
index set, not only its cardinality.  Above \(\finc\), attempts to combine
the source's maximal Bernstein coding with dominated-convergence failure
meet a genuine cardinal-arithmetic bottleneck.

\section{Human-review recommendation}

The packet should receive mathematical review as a candidate substantial
partial result.  The priority checks are:
\begin{enumerate}[leftmargin=*]
\item the passage from the down-directed tail essential unions to the
positive-measure set \(U\) in Lemma~\ref{lem:dominated};
\item the literal-versus-essential liminf comparison in the Fatou case;
\item the interpretation and cofinal pullback used for Q3.
\end{enumerate}

\begin{thebibliography}{9}
\bibitem{source}
D. L. Rodr\'iguez-Vidanes,
\emph{Lineability on nets and uncountable sequences of functions in measure
theory},
arXiv:2306.03862v1 (2023); published in
Rev. R. Acad. Cienc. Exactas F\'is. Nat. Ser. A Mat. 119 (2025), Article 35,
DOI: 10.1007/s13398-024-01696-8.
\end{thebibliography}

\end{document}
